\newpage
\section{Thảo luận}
Dựa trên các kết quả thực nghiệm từ Nhiệm vụ 2 đến Nhiệm vụ 5, nhóm thực hiện đưa ra các phân tích chi tiết nhằm trả lời các câu hỏi về mối liên hệ giữa Docker và lý thuyết hệ điều hành.

\subsection{Ánh xạ các tính năng của Docker vào cơ chế của hệ điều hành}
Thông qua các bài thực hành, nhóm nhận thấy Docker không tự tạo ra sự ảo hóa mà tận dụng trực tiếp các cơ chế có sẵn trong nhân Linux (Host Kernel). Sự ánh xạ cụ thể như sau:

\begin{itemize}
    \item \textbf{Process Isolation (Cách ly tiến trình) $\leftrightarrow$ Linux Namespaces:}
    Docker sử dụng Namespaces để cung cấp cho mỗi container một ``giao diện'' riêng về hệ thống:
    \begin{itemize}
        \item \textit{PID Namespace:} Giúp tiến trình trong container thấy nó có PID = 1 (init process), trong khi ở Host nó chỉ là một tiến trình bình thường với PID khác. Điều này tạo ra sự ảo hóa về bảng tiến trình.
        \item \textit{Mount Namespace:} Tạo ra một hệ thống tệp tin gốc (rootfs) riêng biệt, giúp container không nhìn thấy toàn bộ file system của Host.
        \item \textit{Network Namespace:} Ảo hóa stack mạng (IP, port, routing table riêng).
    \end{itemize}

    \item \textbf{Resource Management (Quản lý tài nguyên) $\leftrightarrow$ Control Groups (cgroups):}
    Cơ chế lập lịch (Scheduling) và quản lý bộ nhớ của OS được Docker điều phối thông qua cgroups:
    \begin{itemize}
        \item Khi giới hạn CPU (\texttt{--cpus}), cgroups can thiệp vào bộ lập lịch của kernel (scheduler) để giới hạn thời gian sử dụng CPU (quota) của nhóm tiến trình container.
        \item Khi giới hạn Memory (\texttt{-m}), cgroups thiết lập ngưỡng cứng. Nếu vượt quá, cơ chế OOM Killer của OS sẽ được kích hoạt cục bộ cho cgroup đó để kết thúc tiến trình.
    \end{itemize}

    \item \textbf{Storage Efficiency (Hiệu quả lưu trữ) $\leftrightarrow$ Union Filesystem (OverlayFS):}
    Docker ánh xạ cơ chế quản lý file của OS thông qua OverlayFS để thực hiện \textbf{Copy-on-Write (CoW)}. Thay vì sao chép toàn bộ dữ liệu khi khởi chạy, OS chỉ ghi nhận các thay đổi ở lớp trên cùng (upperdir). Cơ chế này tối ưu hóa I/O và dung lượng lưu trữ tương tự như cơ chế \textit{Virtual Memory} chia sẻ các trang nhớ (shared pages) giữa tiến trình cha và con sau lệnh \texttt{fork()}.
\end{itemize}

\subsection{Docker gần với VM (máy ảo) hay Sandbox hơn?}
Dựa trên kiến trúc và hành vi quan sát được, \textbf{Docker gần với một Sandbox (cơ chế cô lập) hơn là một Virtual Machine (VM) truyền thống}, nhưng thuật ngữ chính xác nhất là \textbf{ảo hóa cấp hệ điều hành} (OS-level Virtualization).

\begin{itemize}
    \item \textbf{Tại sao không phải là VM?} VM yêu cầu một Hypervisor và một nhân hệ điều hành (Guest Kernel) riêng biệt chạy trên phần cứng ảo hóa. Trong thí nghiệm Task 2, ta thấy container chia sẻ trực tiếp nhân (Kernel) với Host. Các system call từ container (quan sát qua \texttt{strace} ở Task 5) được gửi trực tiếp xuống Host Kernel mà không qua lớp giả lập phần cứng nào.
    
    \item \textbf{Tại sao giống Sandbox?} Giống như Sandbox, Docker giới hạn những gì tiến trình có thể nhìn thấy và tài nguyên nó có thể dùng (thông qua Namespaces và cgroups). Tuy nhiên, Docker mạnh hơn Sandbox thông thường vì nó cung cấp trọn vẹn một môi trường user-space (có shell, thư viện, package manager riêng) tạo cảm giác như một OS độc lập.
\end{itemize}

\textbf{Kết luận:} Docker là một hình thức cô lập tiến trình nâng cao (Advanced Process Isolation). Nó loại bỏ gánh nặng (overhead) của việc chạy một kernel riêng nhưng vẫn đảm bảo tính tách biệt cần thiết.

\subsection{Lựa chọn Base Image phù hợp dưới góc độ hệ điều hành}
Từ kết quả so sánh ở Nhiệm vụ 5 giữa Ubuntu, Debian Slim và Alpine, nhóm đưa ra bảng phân tích sau:

\begin{table}[H]
    \centering
    \renewcommand{\arraystretch}{1.5}
    \begin{tabularx}{\textwidth}{|l|X|X|}
    \hline
    \textbf{Base Image} & \textbf{Đặc điểm OS} & \textbf{Phù hợp cho (Use Case)} \\
    \hline
    \textbf{Alpine Linux} & 
    \textbf{Ưu điểm:} Kích thước cực nhỏ ($\approx$5MB), bảo mật cao. \newline 
    \textbf{Nhược điểm:} Dùng thư viện \texttt{musl libc} thay vì \texttt{glibc}. Có thể gây lỗi tương thích hoặc thay đổi hành vi system call. & 
    \textbf{Microservices, Cloud-native:} Thích hợp cho ứng dụng Go, Rust hoặc static binary không phụ thuộc dynamic linking của hệ thống. \\
    \hline
    \textbf{Ubuntu / Debian} & 
    \textbf{Ưu điểm:} Dùng \texttt{glibc} chuẩn, đầy đủ utilities, dễ debug. \newline 
    \textbf{Nhược điểm:} Kích thước lớn, chứa nhiều thành phần thừa. & 
    \textbf{Môi trường Dev, AI/ML:} Dành cho các ứng dụng cần thư viện C/C++ native phức tạp hoặc cần môi trường giống OS đầy đủ để debug. \\
    \hline
    \textbf{Debian Slim} & 
    \textbf{Đặc điểm:} Bản rút gọn của Debian, vẫn dùng \texttt{glibc} nhưng bỏ các file docs/man pages. & 
    \textbf{Production (General):} Cân bằng tốt nhất giữa kích thước và tương thích. Tốt cho Java, Python, Node.js. \\
    \hline
    \end{tabularx}
    \caption{So sánh các loại base image dưới góc độ hệ điều hành}
\end{table}

\textbf{Góc nhìn hệ điều hành:} Việc chọn image thực chất là chọn \textbf{môi trường user-space}. Vì kernel luôn là của Host, nên sự khác biệt chỉ nằm ở thư viện chuẩn C (libc) và các công cụ hỗ trợ. Nếu ứng dụng phụ thuộc chặt chẽ vào các hành vi cụ thể của \texttt{glibc}, việc dùng Alpine (\texttt{musl}) sẽ gây rủi ro về độ ổn định dù tiết kiệm tài nguyên.